\section{Earlier-period fixed selection rules}\label{sup:selection}

The earlier-period experiment applies the saved fixed selection rules
to familiar historical data. These origins already have validation and
evaluation roles in the current archive, including actual rolling-fold
scoring. They cannot now be treated as untouched evidence for the expanded
backcast. The experiment remains retrospective, not external or prospective
validation. R0 denotes the fixed minimum-validation-loss rule. R2 is the predesignated main one-standard-error selection rule,
not a result-dependent replacement chosen for its subsequent performance.
The rule design precedes its calculations but follows earlier research results.

R1 adds the ordinary empirical quantile to the finite regression candidates.
R2 compares each candidate's validation losses with the same best-mean
candidate, using the paired-difference Bartlett-HAC standard error. Candidates
within that tolerance qualify, and the fixed rule prefers the strongest
shrinkage, including the empirical report. This paired HAC tolerance is a
time-series adaptation of a one-standard-error preference, not the textbook
formula, a significance test or a noninferiority guarantee. It does not account
for the whole uncertainty from selecting the best candidate on that segment.

R3 applies the same preference to prespecified earlier rolling folds and the
outer validation segment, with origin-equal weighting and maturity gaps.
Its HAC pairs use actual source-index distances rather than compressing gaps.
Without an eligible earlier fold, the saved design makes it identical to R2;
this is not a numerical fallback. Earlier-period results were sealed before
the already-inspected-period calculations that could reuse those origins.
\begingroup
\footnotesize
\setlength{\tabcolsep}{3pt}
\renewcommand{\arraystretch}{1.12}
\begin{longtable}{p{35mm}rrrrr}
\caption{Every earlier-period minimum-loss and one-standard-error comparison. The one-standard-error rule misses the recorded success conditions.}\\
\toprule
Series & \shortstack{Horizon /\\history} & \shortstack{Minimum-loss\\rule (\%)} & \shortstack{One-standard-\\error rule (\%)} & \shortstack{Breaches\\minimum-loss /\\one-standard-error} & \shortstack{Historical\\breaches} \\
\midrule
\endfirsthead
\caption[]{Continued.}\\
\toprule
Series & \shortstack{Horizon /\\history} & \shortstack{Minimum-loss\\rule (\%)} & \shortstack{One-standard-\\error rule (\%)} & \shortstack{Breaches\\minimum-loss /\\one-standard-error} & \shortstack{Historical\\breaches} \\
\midrule
\endhead
\bottomrule
\endfoot
PDPOSGS-B & 6/52 & 11.22 & 7.43 & 10/10 & 7 \\
PDPOSGS-B & 24/52 & 181.39 & 181.39 & 27/27 & 8 \\
PDPOSMBS-TOT & 6/52 & -4.11 & 3.83 & 7/6 & 4 \\
PDPOSMBS-TOT & 24/52 & 48.55 & 26.55 & 14/13 & 10 \\
PDPOSCS-TOT & 6/52 & -11.56 & -11.56 & 0/0 & 0 \\
PDPOSCS-TOT & 24/52 & -4.12 & 1.02 & 4/4 & 2 \\
PDPOSFGS-TOT & 6/52 & -27.32 & -27.32 & 6/6 & 0 \\
PDPOSFGS-TOT & 24/52 & -3.87 & -3.60 & 4/3 & 2 \\
PDPOSTIPS-L2 & 6/52 & -7.06 & -7.06 & 2/2 & 1 \\
PDPOSTIPS-L2 & 24/52 & -23.33 & -1.55 & 4/1 & 1 \\
PDPOSTIPS-G11 & 6/52 & -5.55 & -5.79 & 5/5 & 4 \\
PDPOSTIPS-G11 & 24/52 & 5.67 & -0.87 & 15/15 & 15 \\
PDPOSTIPS-G2 & 6/52 & 15.16 & 15.16 & 8/8 & 0 \\
PDPOSTIPS-G2 & 24/52 & 30.27 & 30.27 & 14/14 & 10 \\
PDPOSTIPS-G6L11 & 6/52 & 283.88 & 28.58 & 17/6 & 0 \\
PDPOSTIPS-G6L11 & 24/52 & -4.03 & -2.70 & 2/1 & 0 \\
PDPOSGSC-L2 & 6/52 & 48.04 & 16.09 & 17/15 & 14 \\
PDPOSGSC-L2 & 24/52 & 40.77 & 40.77 & 34/34 & 30 \\
PDPOSGSC-G2L3 & 6/52 & 1.40 & 11.14 & 12/13 & 8 \\
PDPOSGSC-G2L3 & 24/52 & -11.22 & 15.48 & 15/20 & 15 \\
PDPOSGSC-G3L6 & 6/52 & 2.63 & 2.35 & 11/11 & 10 \\
PDPOSGSC-G3L6 & 24/52 & 15.62 & 6.92 & 29/28 & 26 \\
PDPOSGSC-G6L7 & 6/52 & 19.03 & 19.03 & 6/6 & 1 \\
PDPOSGSC-G6L7 & 24/52 & 2.05 & 1.82 & 20/19 & 19 \\
PDPOSGSC-G7L11 & 6/52 & 19.18 & 12.77 & 13/13 & 10 \\
PDPOSGSC-G7L11 & 24/52 & 8.48 & 2.70 & 29/29 & 27 \\
fx\_hourly & 6/720 & 2.44 & 0.62 & 79/73 & 63 \\
fx\_hourly & 6/2880 & -16.62 & -6.28 & 20/20 & 16 \\
fx\_hourly & 24/720 & -4.36 & -4.36 & 47/47 & 32 \\
fx\_hourly & 24/2880 & -1.42 & -1.40 & 1/1 & 0 \\
\end{longtable}
\endgroup

All three recorded success criteria are marked NOT\_MET.
R2's worst weekly excess is 181.39 percent,
and its weekly median is 5.37 percent.
The long-history FX gain also fails the prewritten retention criterion.
This adverse result remains a robustness finding rather than a contribution
claim. Complete inspected-period and all-rule rows remain linked in
\href{https://github.com/ByronLi0207/financial-mathematics-paper-20261020/blob/revision-v7-final/results/supervisor-20261023/r6/README.md}{the R6 archive}.